\documentclass[%
	a4paper,
	twoside,
	openany,
	dvipsnames
]
{book}
%
%----------------------------------------------------------------------------
%
% Title
%
\title{FPGA Based Embedded Neural Network Object Detector}
%
%----------------------------------------------------------------------------
%
% Author
%
\author{Lukas Baischer}
%
%----------------------------------------------------------------------------
%
% Matrikelnummer/Registrationnumber
%
\newcommand{\registrationnumber}{01426506}
%
%----------------------------------------------------------------------------
%
% Supervisor(s)
%
\newcommand{\supervisor}{Univ.Prof. Dipl.-Ing. Dr.techn. Axel  Jantsch}
%
%----------------------------------------------------------------------------
%
% Submission Date
%
\newcommand{\submissiondate}{27.10.2021}

%
% ----------------------------------------------------------------------------
%
% Uncomment only one of the following depending on your type of document:
%
%\newcommand{\doctype}{REPORT}
%\newcommand{\doctype}{DISSERTATION}
\newcommand{\doctype}{MASTERTHESIS}
%\newcommand{\doctype}{BACHELORTHESIS}
%
%-----------------------------------------------------------------------------
%
% Use the 'Libertine' font type
%
\usepackage{libertine}
\usepackage[T1]{fontenc}
\usepackage[utf8]{inputenc}
\usepackage[english,ngerman]{babel}
%
%----------------------------------------------------------------------------
%
% Set page margins
%
\usepackage{geometry}
\geometry{%
	left   = 2cm,
	right  = 2cm,
	top    = 2cm,
	bottom = 2cm
}
%
%----------------------------------------------------------------------------
%
% Set line spacing
%
\usepackage{setspace}
\setstretch{1.5}
%
%----------------------------------------------------------------------------
%
% Set paragraph: No indentation, but include an empty line
%
\usepackage[parfill]{parskip}
%
%----------------------------------------------------------------------------
%
% Settings for hyperlinks
%
\usepackage{hyperref}
\hypersetup{%
	colorlinks = false,
	allcolors  = blue,
}
%
%----------------------------------------------------------------------------
%
% Use graphics
%
\usepackage{graphicx}
\usepackage{subcaption }
%
%----------------------------------------------------------------------------
%
% Use colors
%
\usepackage{xcolor}
\usepackage{colortbl}
%
%----------------------------------------------------------------------------
%
% Define a TODO and a DONE command
%
\newcommand{\todo}[1]{\textcolor{red}{#1}}
\newcommand{\done}[1]{}
%
%
\usepackage[hang]{footmisc}
\renewcommand{\hangfootparindent}{2em}
\renewcommand{\hangfootparskip}{2em}
\renewcommand{\footnotemargin}{0.00001pt}
\renewcommand{\footnotelayout}{\hspace{2em}}
%
%
%----------------------------------------------------------------------------
%
% Set line spacing for all figure environments
% From: https://tex.stackexchange.com/a/166458
%
\let\svfigure\figure
\let\svendfigure\endfigure
\renewenvironment{figure}[1][tb]{\svfigure[#1]\setstretch{1}}
{\svendfigure}
%
%----------------------------------------------------------------------------
%
% Use AMS math fonts
%
\usepackage{amsfonts}
\usepackage[sans]{dsfont}
%
%----------------------------------------------------------------------------
%
% Use multiple figures in one float
%
\usepackage{subcaption}
%
%----------------------------------------------------------------------------
%
% Use dummy text
%
\usepackage{lipsum}
%
%----------------------------------------------------------------------------
%
% Use extended list environments (e.g., 'inparaenum')
%
\usepackage{paralist}
%
%----------------------------------------------------------------------------
%
% Use listings
%
\usepackage{listings}
%
%----------------------------------------------------------------------------
%
% Typeset pseudo code
%
\usepackage{syntax}
%
%----------------------------------------------------------------------------
%
% More options for boxes
%
\usepackage{realboxes}
%
% Command for vertical text in tabulars
%
\newcommand*\rot{\rotatebox{90}}
%
%----------------------------------------------------------------------------
%
% Package for logos (e.g., the BibTeX logo)
%
\usepackage{dtk-logos}
%
%----------------------------------------------------------------------------
%
% Use \textsubscript
%
\usepackage{fixltx2e}
%
%----------------------------------------------------------------------------
%
% More options for tabulars
%
\usepackage{array}
%
%----------------------------------------------------------------------------
%
% Use appendices
%
\usepackage[titletoc]{appendix}
%
%----------------------------------------------------------------------------
%
% Macros to use title and author information in the document
%
\makeatletter
\let\thetitle\@title
\let\theauthor\@author
\makeatother
%
%----------------------------------------------------------------------------
%
% Do not show page numbers on empty pages
%
\usepackage{emptypage}
%
%---------------------------------------------------------------------------
%
% Set spacing between items (looks really weird otherwise)
\usepackage{enumitem}
\setitemize{itemsep=-8pt}
\setenumerate{itemsep=-8pt}
%
%----------------------------------------------------------------------------
%
% Use algorithm and algpseudocode to create pseudo code algorithms 
%
\usepackage{algorithm} 
\usepackage{algpseudocode} 
%----------------------------------------------------------------------------
%
% Use pythonhighlight for python syntax highlighting see https://github.com/olivierverdier/python-latex-highlighting for more information
%
\usepackage{pythonhighlight}
%
%----------------------------------------------------------------------------
%
% Use svg package for plotting inkscape graphics 
%
\usepackage{svg}
%
%----------------------------------------------------------------------------
% 
% Use float to be able to force figure position using [H]
%
\usepackage{float}
%
%----------------------------------------------------------------------------
% 
% Use pdfpages to be able to add pdf to appendix and pdflscape for landcape pdfs 
%
\usepackage{pdfpages}
\usepackage{pdflscape}
%
%----------------------------------------------------------------------------
% 
% Use glossaries to use acronyms and symbols 
%
\usepackage[acronym,nonumberlist,nowarn,style=long]{glossaries}
\usepackage{makeidx}
\usepackage{csquotes}
%
%----------------------------------------------------------------------------
% 
% For Si units 
%
\usepackage{siunitx}
%
%----------------------------------------------------------------------------
% 
% Box for for research question 
%
\usepackage{fbox}
%
%----------------------------------------------------------------------------
% 
% For nicer tables 
%
\usepackage{booktabs}
\usepackage{tabularx}
\newcolumntype{C}{>{\centering\arraybackslash}X} 
\usepackage{multirow}
%
%----------------------------------------------------------------------------
%
% For section comments 
%
\usepackage{comment}
%
%----------------------------------------------------------------------------
% 
% cite in order to sort citation numbers  
%
\usepackage{cite}
%
%----------------------------------------------------------------------------
%
% Use the cleverref package -- Load this package as the very last!
%
\usepackage{nameref}
%
%----------------------------------------------------------------------------
%
% Use the cleverref package -- Load this package as the very last!
%
\usepackage{cleveref}
%
%----------------------------------------------------------------------------
%
% Include listings for Code highlighting 
%


\definecolor{mygreen}{RGB}{0,102,0} % color values Red, Green, Blue
\definecolor{mylilas}{RGB}{170,55,241}
\definecolor{myred}{RGB}{94,32,40}

%----------------------------------------------------------------------------
% Matlab
\lstset{language=Matlab,%
    %basicstyle=\color{red},
    frame = single,
    basicstyle=\tiny,
    breaklines=true,%
    morekeywords={matlab2tikz},
    keywordstyle=\color{blue},%
    morekeywords=[2]{1}, keywordstyle=[2]{\color{black}},
    identifierstyle=\color{black},%
    stringstyle=\color{mylilas},
    commentstyle=\color{mygreen},%
    showstringspaces=false,%without this there will be a symbol in the places where there is a space
    numbers=left,%
    numberstyle={\tiny \color{black}},% size of the numbers
    numbersep=9pt, % this defines how far the numbers are from the text
    emph=[1]{for,end,break},emphstyle=[1]\color{red}, %some words to emphasise
    %emph=[2]{word1,word2}, emphstyle=[2]{style},    
}
%----------------------------------------------------------------------------
% C
\lstdefinestyle{CStyle}{
frame = single,
    backgroundcolor=\color{white},   
    commentstyle=\color{mygreen},
    keywordstyle=\color{blue},
    numberstyle=\tiny\color{black},
    stringstyle=\color{mylilas},
    basicstyle=\tiny,
    breakatwhitespace=false,         
    breaklines=true,                 
    captionpos=b,                    
    keepspaces=true,                 
    numbers=left,                    
    numbersep=9pt,                  
    showspaces=false,                
    showstringspaces=false,
    showtabs=false,                  
    tabsize=2,
    emph=[1]{for,end,break},emphstyle=[1]\color{blue}, %some words to emphasise
    language=C
}
%----------------------------------------------------------------------------
% VHDL
\lstdefinelanguage{VHDL1}{
    breaklines=true,%
    numberstyle={\tiny \color{black}},% size of the numbers
    numbersep=9pt, % this defines how far the numbers are from the text
   morekeywords={
     library,use,all,entity,is,port,in,out,end,architecture,of,
     begin,and,else,if,elsif,process,signal,end,then
   },
    emph=[1]{std_logic,std_logic_vector,STD_LOGIC,STD_LOGIC_VECTOR,rising_edge},emphstyle=[1]\color{myred},
   morecomment=[l]--
}

\lstdefinestyle{vhdl}{
    frame = single, 
    backgroundcolor=\color{white},   
    numberstyle=\tiny\color{black},
    stringstyle=\color{mylilas},
    basicstyle=\tiny,
    breakatwhitespace=false,         
    breaklines=true,                 
    captionpos=b,                    
    keepspaces=true,                 
    numbers=left,                    
    numbersep=9pt,                  
    showspaces=false,                
    showstringspaces=false,
    showtabs=false,                  
    tabsize=4,
   language     = VHDL1,
   keywordstyle = \color{blue}\bfseries,
   commentstyle = \color{mygreen}
}

%
%----------------------------------------------------------------------------
%
%----------------------------------------------------------------------------
%
% GLOSSARIES SETTINGS
%:::::::::::::::::::::::::::::::::::::::::::::::::::::::::::::::::::::::::::::::::::::::::::::::::::::::::::::
\setlength{\glsdescwidth}{1.0\textwidth}		% left aligned
\renewcommand*{\glspostdescription}{}				% Remove the dot at the end of glossary descriptions
\renewcommand*{\glsgroupskip}{}							% Remove vertical space between index groups
%:::::::::::::::::::::::::::::::::::::::::::::::::::::::::::::::::::::::::::::::::::::::::::::::::::::::::::::
%
%
% SYMBOLS, ACRONYMS AND INDEX
%:::::::::::::::::::::::::::::::::::::::::::::::::::::::::::::::::::::::::::::::::::::::::::::::::::::::::::::
\newglossary[slg]{symbolslist}{syi}{syg}{Symbols}
%
\newacronym{SNR}{SNR}{signal to noise ration}
\newacronym{FPGA}{FPGA}{field programmable gate array}
\newacronym{IP}{IP}{intellectual property}
\newacronym{ILA}{ILA}{integrated logic analyzer}
\newacronym{ADC}{ADC}{analog to digital converter}
\newacronym{CNN}{CNN}{convolutional neural network}
\newacronym{DMA}{DMA}{direct memory access}
\newacronym{PS}{PS}{processing system}
\newacronym{PL}{PL}{programmable logic}
\newacronym{BRAM}{BRAM}{block-RAM}
\newacronym{PE}{PE}{processing element}
\newacronym{PU}{PU}{processing unit}
\newacronym{LE}{LE}{logic element}
\newacronym{MAC}{MAC}{multiply and accumulate}
\newacronym{ASIC}{ASIC}{application specific integrated circuit}
\newacronym{GPU}{GPU}{graphics processing unit}
\newacronym{CPU}{CPU}{central processing unit}
\newacronym{RNN}{RNN}{recurrent neural network}
\newacronym{FC}{FC}{fully connected}
\newacronym{DNN}{DNN}{deep neural network}
\newacronym{ANN}{ANN}{artificial neural network}
\newacronym{SNN}{SNN}{spiking neural network}
\newacronym{BNN}{BNN}{binarized neural network}
\newacronym{TNN}{TNN}{ternary neural network}
\newacronym{ALU}{ALU}{arithmetic and logic unit}
\newacronym{AI}{AI}{artificial intelligence}
\newacronym{LUT}{LUT}{lookup table}
\newacronym{SOC}{SoC}{system on chip}
\newacronym{GEMM}{GEMM}{general matrix multiplication}
\newacronym{SIMD}{SIMD}{single instruction multiple data}
\newacronym{SIMT}{SIMT}{single instruction multiple threads}
\newacronym{OpenCL}{OpenCL}{open computer language}
\newacronym{FIR}{FIR}{finite impulse response}
\newacronym{TPU}{TPU}{tensor processing unit}
\newacronym{AOI}{AOI}{automated optical inspection}
\newacronym{MOI}{MOI}{manual optical inspection}
\newacronym{PCB}{PCB}{printed circuit board}
\newacronym{PCBA}{PCBA}{printed circuit board assembly}
\newacronym{MV}{MV}{machine vision}
\newacronym{DSP}{DSP}{digital signal processor}
\newacronym{IC}{IC}{integrated circuit}
\newacronym{HLS}{HLS}{high level synthesis}
\newacronym{RTL}{RTL}{register transfer level}
\newacronym{FFT}{FFT}{fast Fourier transformation}
\newacronym{FF}{FF}{flip flop}
\newacronym{SP}{SP}{soft processor}
\newacronym{CLB}{CLB}{configurable logic block}
\newacronym{Relu}{ReLU}{rectified linear unit}
\newacronym{NoC}{NoC}{network on chip}
\newacronym{SSD}{SSD}{single shot multibox detector}
\newacronym{YOLO}{YOLO}{you only look ones}
\newacronym{LSTM}{LSTM}{long short term memory}
\newacronym{GRU}{GRU}{gated recurrent unit}
\newacronym{v4l2}{v4l2}{video for linux 2}
\newacronym{mlp}{MLP}{multilayer perceptron}
\newacronym{NFU}{NFU}{neural function unit}
\newacronym{IOU}{IoU}{intersection over union}
\newacronym{mAP}{mAP}{mean average precision}
\newacronym{nms}{NMS}{none maximum surpression}
\newglossaryentry{symb:lambda}{name={$\lambda$},description={Wavelength},type=symbolslist}
\newglossaryentry{symb:tau}{name={$\tau$},description={delay},type=symbolslist}
\newglossaryentry{symb:epsilon}{name={$\epsilon$},description={Measurement error},type=symbolslist}
\newglossaryentry{symb:frequency}{name={$f$},description={Frequency},type=symbolslist}
\newglossaryentry{symb:Gamma}{name={$\Gamma$},description={Reflection coefficient},type=symbolslist}
\newglossaryentry{symb:V}{name={$V$},description={Voltage},type=symbolslist}
\newglossaryentry{symb:I}{name={$I$},description={Current},type=symbolslist}
\newglossaryentry{symb:Z}{name={$Z$},description={Impedance},type=symbolslist}
\newglossaryentry{symb:G}{name={$G$},description={Gain},type=symbolslist}
\newglossaryentry{symb:R}{name={$R$},description={Resistance},type=symbolslist}
\newglossaryentry{symb:C}{name={$C$},description={Capacitance},type=symbolslist}
\newglossaryentry{symb:T}{name={$T$},description={Temperature},type=symbolslist}
\newglossaryentry{symb:Jpp}{name={$J_p_p$},description={Peak to peak jitter},type=symbolslist}
\newglossaryentry{symb:Jstd}{name={$J_\sigma$},description={Standard deviation of jitter},type=symbolslist}
\newglossaryentry{symb:Pre}{name={$\epsilon_\sigma$},description={Precision},type=symbolslist}
\newglossaryentry{symb:maxErr}{name={$\epsilon_m_a_x$},description={Maximum simulation error},type=symbolslist}
\newglossaryentry{symb:sps}{name={Sps},description={Samples per second},type=symbolslist}
\newglossaryentry{fps}{name={fps},description={frames per second},type=symbolslist}
\newglossaryentry{GOPS}{name={GOPS},description={giga operations per second},type=symbolslist}
\newglossaryentry{FLOPS}{name={FLOP/s},description={floating point operations per second},type=symbolslist}

%
\makeglossaries
\makeindex


% Document body
%
\begin{document}
\selectlanguage{english}
%
%----------------------------------------------------------------------------
%
% Select title page
%
\ifthenelse{\equal{\doctype}{REPORT}}{% Report title page

\begin{titlepage}

	\begin{center}

	\includegraphics[height=2cm]{fig/logo-tu-bw.png}%
	\hfill{}%
	\includegraphics[height=2cm]{fig/logo-ict.png}%
	

	\vspace{5em}

	{\Huge Report}
	\vspace{3em}

	{\large submitted by}

	\vspace{3em}

	{\huge \theauthor}

	{\large Registration Number \registrationnumber}

	\vspace{5em}

	{\Huge \thetitle}

	\vspace{3em}
	{\large SoC Advanced Course  }

	\end{center}

	\vspace{3em}

	\large
	\begin{tabular}{m{.5\textwidth}m{.5\textwidth}}
	Vienna, Austria, \submissiondate & \\
	\end{tabular}

	\vspace{2em}

	\begin{tabular}{m{.5\textwidth}m{.5\textwidth}}
	Study code:     & 066 504 \\
	Field of study: & Electrical Engineering \\
	\end{tabular}

	\vspace{2em}

	\begin{tabular}{m{.5\textwidth}m{.5\textwidth}}
	Supervisor:    & \supervisor \\
	\ifdef{\cosupervisor}{
	Co-Supervisor: & \cosupervisor \\
	}{}
	\end{tabular}

\end{titlepage}


Copyright (C) 2018 \theauthor

If you find this work useful, please cite it using the following \BibTeX{ } entry:

\vspace{1em}

\begin{lstlisting}[%
	breaklines = true,%
	basicstyle = \ttfamily\footnotesize,%
 escapeinside={(*@}{@*)},
	keepspaces = true,
	frame      = single,%
]
@TechReport{AuthorLastName2021,
  author = 	 {(*@\theauthor@*)},
  title = 	 {(*@\thetitle@*)},
  institution =  {TU Wien},
  year = 	 {2021},
  address = 	 {Gusshausstrasse 27--29 / 384, 1040 Wien, Austria},
  month = 	 {November}
}
\end{lstlisting}

\vspace{3em}
Contact us:

\href{E-mail address}{author@institution.domain}

\vfill

\includegraphics[height=1.5cm]{fig/cc-large.png}
\includegraphics[height=1.5cm]{fig/by-large.png}


This report is licensed under the following license:
Attribution 4.0 International (CC BY 4.0)

\vspace{3em}

You are free to:

\begin{enumerate}
   \item Share --- Copy and redistribute the material in any medium or format
   \item Adapt --- Remix, transform, and build upon the material for any purpose,
			even commercially.
\end{enumerate}

This license is acceptable for Free Cultural Works.

The licensor cannot revoke these freedoms as long as you follow the license terms.

The entire license text is available at:
\href{https://creativecommons.org/licenses/by/4.0/legalcode}
	{https://creativecommons.org/licenses/by/4.0/legalcode}
	
\pagebreak

%% End of REPORT title page

%%% Local Variables:
%%% mode: latex
%%% TeX-master: "thesis"
%%% End:
}{}
\ifthenelse{\equal{\doctype}{DISSERTATION}}{%----------------------------------------------------------------------------
%
%  Dissertation
%
\begin{titlepage}

	\begin{center}

	\includegraphics[height=2cm]{fig/logo-tu-bw.png}%
	\hfill{}%
	\includegraphics[height=2cm]{fig/logo-ict.png}%
	

	\vspace{5em}

	{\Huge Doctoral Thesis}
	\vspace{3em}

	{\large submitted by}

	\vspace{3em}

	{\huge \theauthor}

	{\large Registration Number \registrationnumber}

	\vspace{5em}

	{\Huge \thetitle}

	\vspace{3em}
	{\large In partial fulfillment of the requirements for the degree of}

	\vspace{3em}

	{\Large Doctor of Philosophy (PhD)}

	\end{center}

	\vspace{3em}

	\large
	\begin{tabular}{m{.5\textwidth}m{.5\textwidth}}
	Vienna, Austria, \submissiondate & \\
	\end{tabular}

	\vspace{2em}

	\begin{tabular}{m{.5\textwidth}m{.5\textwidth}}
	Study code:     & E\,786\,710 \\
	Field of study: & Electrical Engineering \\
	\end{tabular}

	\vspace{2em}

	\begin{tabular}{m{.5\textwidth}m{.5\textwidth}}
	Supervisor:    & \supervisor \\
	Co-Supervisor: & \cosupervisor \\
	\end{tabular}

\end{titlepage}


Copyright (C) 2018 \theauthor

If you find this work useful, please cite it using the following \BibTeX{ } entry:

\vspace{1em}

\begin{lstlisting}[%
	breaklines = true,%
	basicstyle = \ttfamily\footnotesize,%
  escapeinside={(*@}{@*)},
	keepspaces = true,
	frame      = single,%
]
@Thesis{AuthorLastname2018,
  type        = {Doctoral Thesis},
  author      = {(*@\theauthor@*)},
  title       = {(*@\thetitle@*)},
  school      = {Vienna University of Technology (TU Wien)},
  year        = {2018},
  address     = {Gusshausstrasse 27--29 / 384, 1040 Wien},
  month       = {November},
}
\end{lstlisting}

\vspace{3em}
Contact us:

\href{E-mail address}{author@institution.domain}

\vfill

\includegraphics[height=1.5cm]{fig/cc-large.png}
\includegraphics[height=1.5cm]{fig/by-large.png}


This dissertation is licensed under the following license:
Attribution 4.0 International (CC BY 4.0)

\vspace{3em}

You are free to:

\begin{enumerate}
    \item Share --- Copy and redistribute the material in any medium or format
    \item Adapt --- Remix, transform, and build upon the material for any purpose,
			even commercially.
\end{enumerate}

This license is acceptable for Free Cultural Works.

The licensor cannot revoke these freedoms as long as you follow the license terms.

The entire license text is available at:
\href{https://creativecommons.org/licenses/by/4.0/legalcode}
	{https://creativecommons.org/licenses/by/4.0/legalcode}

%%% Local Variables:
%%% mode: latex
%%% TeX-master: "thesis"
%%% End:
}{}
\ifthenelse{\equal{\doctype}{MASTERTHESIS}}{
%----------------------------------------------------------------------------
%
% Master's Thesis
%
\begin{titlepage}

	\begin{center}

	\includegraphics[height=2cm]{fig/logo-tu-bw.png}%
	\hfill{}%
	\includegraphics[height=2cm]{fig/logo-ict.png}%


	\vspace{5em}

	{\Huge Master's Thesis}
	\vspace{2em}

	{\large submitted by}

	\vspace{3em}

	{\huge \theauthor}

	{\large Registration Number \registrationnumber}

	\vspace{3em}

	{\Huge \thetitle}

	\vspace{3em}
	{\large In partial fulfillment of the requirements for the degree of}

	\vspace{3em}

	{\Large Diplom-Ingenieur (Dipl.-Ing.)}

	\end{center}

	\vspace{3em}

	\large
	\begin{tabular}{m{.5\textwidth}m{.5\textwidth}}
	Vienna, Austria, \submissiondate & \\
	\end{tabular}

	\vspace{2em}

	\begin{tabular}{m{.5\textwidth}m{.5\textwidth}}
	Study code:     & 066 504 \\
	Field of study: & Embedded Systems \\
	\end{tabular}

	\vspace{2em}

	\begin{tabular}{m{.5\textwidth}m{.5\textwidth}}
	Supervisor:    & \supervisor \\
	%Co-Supervisor: & \cosupervisor \\
	\end{tabular}

\end{titlepage}



Copyright (C) 2021 \theauthor

If you find this work useful, please cite it using the following \BibTeX{ } entry:

\vspace{1em}

\begin{lstlisting}[%
	breaklines = true,%
	basicstyle = \ttfamily\footnotesize,%
 escapeinside={(*@}{@*)},
	keepspaces = true,
	frame      = single,%
]
@Thesis{baischer2021,
 type        = {Master's Thesis},
 author      = {(*@\theauthor@*)},
 title       = {(*@\thetitle@*)},
 school      = {Vienna University of Technology (TU Wien)},
 year        = {2021},
 address     = {Gusshausstrasse 27--29 / 384, 1040 Wien},
 month       = {October},
}
\end{lstlisting}

\vspace{3em}
Contact me:

\href{E-mail address}{lukas.baischer@student.tuwien.ac.at}

\vfill

\includegraphics[height=1.5cm]{fig/cc-large.png}
\includegraphics[height=1.5cm]{fig/by-large.png}


This thesis is licensed under the following license:
Attribution 4.0 International (CC BY 4.0)

\vspace{3em}

You are free to:

\begin{enumerate}
   \item Share --- Copy and redistribute the material in any medium or format
   \item Adapt --- Remix, transform, and build upon the material for any purpose,
			even commercially.
\end{enumerate}

This license is acceptable for Free Cultural Works.

The licensor cannot revoke these freedoms as long as you follow the license terms.

The entire license text is available at:
\href{https://creativecommons.org/licenses/by/4.0/legalcode}
	{https://creativecommons.org/licenses/by/4.0/legalcode}


%%% Local Variables:
%%% mode: latex
%%% TeX-master: "thesis"
%%% End:
}{}
\ifthenelse{\equal{\doctype}{BACHELORTHESIS}}{%----------------------------------------------------------------------------
%
% Bachelor's Thesis
%
\begin{titlepage}

	\begin{center}

	\includegraphics[height=2cm]{fig/logo-tu-bw.png}%
	\hfill{}%
	\includegraphics[height=2cm]{fig/logo-ict.png}%


	\vspace{5em}

	{\Huge Bachelor's Thesis}
	\vspace{2em}

	{\large submitted by}

	\vspace{3em}

	{\huge \theauthor}

	{\large Registration Number \registrationnumber}

	\vspace{3em}

	{\Huge \thetitle}

	\vspace{3em}
	{\large In partial fulfillment of the requirements for the degree of}

	\vspace{3em}

	{\Large Bachelor of Science (BSc)}

	\end{center}

	\vspace{3em}

	\large
	\begin{tabular}{m{.5\textwidth}m{.5\textwidth}}
	Vienna, Austria, \submissiondate & \\
	\end{tabular}

	\vspace{2em}

	\begin{tabular}{m{.5\textwidth}m{.5\textwidth}}
	Study code:     & 033 235 \\
	Field of study: & Electrical Engineering \\
	\end{tabular}

	\vspace{2em}

	\begin{tabular}{m{.5\textwidth}m{.5\textwidth}}
	Supervisor:    & \supervisor \\
	Co-Supervisor: & \cosupervisor \\
	\end{tabular}

\end{titlepage}



Copyright (C) 2018 \theauthor

If you find this work useful, please cite it using the following \BibTeX{ } entry:

\vspace{1em}

\begin{lstlisting}[%
	breaklines = true,%
	basicstyle = \ttfamily\footnotesize,%
 escapeinside={(*@}{@*)},
	keepspaces = true,
	frame      = single,%
]
@Thesis{AuthorLastName2018,
 type        = {Bachelor's Thesis},
 author      = {(*@\theauthor@*)},
 title       = {(*@\thetitle@*)},
 school      = {Vienna University of Technology (TU Wien)},
 year        = {2018},
 address     = {Gusshausstrasse 27--29 / 384, 1040 Wien},
 month       = {November},
}
\end{lstlisting}

\vspace{3em}
Contact me:

\href{E-mail address}{author@institution.domain}

\vfill

\includegraphics[height=1.5cm]{fig/cc-large.png}
\includegraphics[height=1.5cm]{fig/by-large.png}


This thesis is licensed under the following license:
Attribution 4.0 International (CC BY 4.0)

\vspace{3em}

You are free to:

\begin{enumerate}
   \item Share --- Copy and redistribute the material in any medium or format
   \item Adapt --- Remix, transform, and build upon the material for any purpose,
			even commercially.
\end{enumerate}

This license is acceptable for Free Cultural Works.

The licensor cannot revoke these freedoms as long as you follow the license terms.

The entire license text is available at:
\href{https://creativecommons.org/licenses/by/4.0/legalcode}
	{https://creativecommons.org/licenses/by/4.0/legalcode}


%%% Local Variables:
%%% mode: latex
%%% TeX-master: "thesis"
%%% End:
}{}

\pagebreak
%
%-----------------------------------------------------------------------------
% Acknowledgements, Abstract, Kurzfassung
% This is not used for reports
%
\chapter*{Acknowledgements}
\input{../sources/tex/front-matter/Acknowledgements.tex}
\chapter*{Abstract}
\addcontentsline{toc}{chapter}{Abstract}
\input{../sources/tex/front-matter/Abstract.tex}
\chapter*{Kurzfassung}
\selectlanguage{german}
\addcontentsline{toc}{chapter}{Kurzfassung}
\input{../sources/tex/front-matter/Kurzfassung.tex}
\selectlanguage{english}


\tableofcontents
\listoffigures 
\listoftables
\printglossary[type=symbolslist]							% List of symbols
\printglossary[type=\acronymtype]

%
%----------------------------------------------------------------------------
%############################################################################
% Put your text body here!
\input{../sources/tex/introduction/introduction.tex}
\input{../sources/tex/state-of-the-art/neural-networks.tex}	
\input{../sources/tex/state-of-the-art/object-detectors.tex}	
\input{../sources/tex/state-of-the-art/hardware-accelerator.tex}
\input{../sources/tex/design/hardware-accelerator-design.tex}
\input{../sources/tex/design/software.tex}
\input{../sources/tex/results.tex}
\input{conclusion.tex}
\input{outlook.tex}
%############################################################################
%----------------------------------------------------------------------------
%
%
%----------------------------------------------------------------------------
%
% Bibliography
%
\backmatter
%\printbibliography
\bibliographystyle{../sources/classes-styles/IEEEtran}
\bibliography{../sources/bibliography/bibliography}
% uncomment the following line if you want the bibliography be shown in table of contents
%\addcontentsline{toc}{chapter}{Bibliography}

%
%----------------------------------------------------------------------------
%
% Declaration of Authorship
% This is not used for reports
\ifthenelse{\equal{\doctype}{REPORT}}{}{
\chapter*{Erklärung zur Verfassung der Arbeit}
Hiermit erkläre ich, dass die vorliegende Arbeit gemäß dem Code of Conduct –
Regeln zur Sicherung guter wissenschaftlicher Praxis (in der aktuellen Fassung
des jeweiligen Mitteilungsblattes der TU Wien), insbesondere ohne unzulässige
Hilfe Dritter und ohne Benutzung anderer als der angegebenen Hilfsmittel,
angefertigt wurde. Die aus anderen Quellen direkt oder indirekt übernommenen
Daten und Konzepte sind unter Angabe der Quelle gekennzeichnet. Die Arbeit
wurde bisher weder im In– noch im Ausland in gleicher oder in ähnlicher Form
in anderen Prüfungsverfahren vorgelegt.

\vspace{3cm}

\thispagestyle{empty}
\begin{center}
\begin{tabular}{@{}p{3.5in}p{2.5in}@{}}
  Vienna, Austria \submissiondate & \hrulefill \\
  & \centering \theauthor  \\
\end{tabular}
\end{center}
}
% Appendix
%
\appendix
\input{../sources/tex/appendix/appendixA}

\end{document}
%
%----------------------------------------------------------------------------

%%% Local Variables:
%%% mode: latex
%%% TeX-master: t
%%% End:
